\begin{tikzpicture}[
  every node/.style={font=\footnotesize\sffamily},
  participant/.style={draw=TFGBlue, fill=TFGLightBlue, rounded corners=2pt,
    minimum width=2.25cm, minimum height=0.65cm, align=center},
  lifeline/.style={dashed, draw=TFGMuted},
  message/.style={-{Latex[length=2mm]}, semithick, draw=TFGAccent},
  return/.style={-{Latex[length=2mm]}, dashed, draw=TFGMuted},
  message label/.style={fill=white, inner xsep=2pt, inner ysep=1pt,
    text=black, align=center}
]
  \node[participant] (user) at (0,0) {Miembro};
  \node[participant] (web) at (4,0) {Aplicación web};
  \node[participant] (api) at (8,0) {Gestor de\\oportunidades};
  \node[participant] (workflow) at (12,0) {Gestor de\\workflows};
  \node[participant] (db) at (16,0) {Persistencia};
  \foreach \x in {0,4,8,12,16} \draw[lifeline] (\x,-0.35) -- (\x,-7.15);
  \draw[message] (0,-0.8) -- node[message label, above]{selecciona workflow} (4,-0.8);
  \draw[message] (4,-1.6) -- node[message label, above]{solicita la asignación} (8,-1.6);
  \draw[message] (8,-2.4) -- node[message label, above]{consulta la definición activa} (12,-2.4);
  \draw[return] (12,-3.2) -- node[message label, above]{definición, pasos y acciones} (8,-3.2);
  \draw[message] (8,-4) -- node[message label, above]{crea ejecución e instancias} (16,-4);
  \draw[return] (16,-4.8) -- node[message label, above]{ejecución inicial registrada} (8,-4.8);
  \draw[return] (8,-5.6) -- node[message label, above]{asignación confirmada} (4,-5.6);
  \draw[return] (4,-6.4) -- node[message label, above]{muestra el workflow asignado} (0,-6.4);
\end{tikzpicture}
